% ==========================================================================
% GENERATED by scripts/split_paper.py. Do not edit this file.
%
% Source: imported_papers/grokking/iclr2027_conference_v2.tex
% Part 7: Experimental details
%
% Edits belong upstream, in the published paper. To change what the
% chapter includes, change imports/grokking.toml and re-run
% the splitter.
% ==========================================================================
\section{Experimental details}
\label{app:repro}

\paragraph{Statistical reporting.} Results are means across three random seeds, with standard deviations in parentheses, unless stated otherwise. Conditions that are compared with each other use the same seeds and data splits. Independent experiments have their own reference runs, so reference values vary slightly between tables. For comparisons between credit and circuit participation we take the absolute difference within each seed and then average.

\paragraph{Networks and training.} Training is full batch. Hidden weights use Kaiming normal initialization \citep{he2015delving}. Adam uses learning rate $10^{-3}$, $\beta_1 = 0.9$, $\beta_2 = 0.999$, $\epsilon = 10^{-8}$ and coupled weight decay, in which the decay term is added to the gradient before the moment estimates; biases are trained and decayed like the weights. For the learning signals passed to Adam, the gradient norm is clipped to $1.0$ after any rescaling and before the optimizer step, and Adam adds the decay term to the clipped gradient. The optimizers compared in Section~\ref{sec:law} apply no clipping. In the multilayer perceptrons, credit equalization is applied to the incoming-weight gradients of the first hidden layer (bias gradients are not rescaled), and circuit participation and neuron deletion are measured on the outgoing weights of the same neurons. The transformer experiments apply it to every linear map in the canonical configuration and to the incoming weights of the MLP in the wide configuration (Appendix~\ref{app:transformer}). Seeds start at $0$, so three-seed experiments use seeds $0$ to $2$. Each run uses one GPU. Runs continue past grokking until participation has stabilized, which takes $12{,}000$ to $70{,}000$ epochs depending on task, modulus and architecture; compared conditions use the same budget, and the experiments with other budgets state them.

\paragraph{Modular-arithmetic data.} We hold out $40\%$ of the $p^2$ operand pairs at random and train on the rest, with the same split procedure for all three operations. Multiplication includes the pairs with a zero operand, identically across rules. We test $p \in \{31, 47, 59\}$ and use $p = 31$ unless stated otherwise.

\paragraph{Optimizer settings.} The comparison of Section~\ref{sec:law} runs every optimizer at the reference learning rate $10^{-3}$ and weight decay $10^{-3}$ for $20{,}000$ epochs. LARS uses fifty times this learning rate. Nero uses its default learning rate $10^{-2}$ and no weight decay, since its norm constraint fixes the weight scale; it normalizes each neuron's gradient by a running estimate of its norm and holds each neuron's incoming-weight norm fixed after its first update. Lion adds weight decay to the gradient before the signed momentum update. Muon uses five cubic Newton--Schulz iterations, ordinary momentum and coupled weight decay on matrices. LARS uses a layerwise norm ratio without momentum. LAMB uses uncorrected moments and caps the trust ratio at $10$. Adafactor uses fixed-rate factored second moments without update clipping, Shampoo uses exponential second moments, and RMSProp uses second-moment decay $0.99$ without momentum. Evaluated every $50$ epochs over three seeds, every optimizer fits the training set first and exceeds test accuracy $0.95$ at least $500$ epochs later, so all of them grok under the criterion of Section~\ref{sec:setup}: the delays are $500$ to $650$ epochs for LARS, $650$ to $700$ for Lion, $1200$ to $1350$ for RMSProp and Adafactor, $2150$ to $2400$ for LAMB, $7450$ to $7800$ for Muon, $12{,}200$ to $13{,}100$ for Shampoo and $4450$ to $7000$ for Nero, against $1300$ to $1550$ for the reference.

\paragraph{Deep feedback control.} DFC uses a controller during training that moves the network output toward a target. The output $\hat{\vy}$ is the linear readout, that is, the logits, and the label $\vy$ is one-hot. The output error is $\ve = \vy^* - \hat{\vy}$, where the target $\vy^* = \hat{\vy} + 2\kappa(\vy - \hat{\vy})$ lies between the current output and the label at a distance set by the control strength $\kappa$; it is one step on the squared error to the label, as in the reference implementation. A leaky proportional-integral controller turns the error into a control signal $\vu$, which is fed back into every layer $i$ through the transpose of the Jacobian $\mJ_i$ from that layer to the output:
\begin{align}
\vu_{\mathrm{int}} &\leftarrow \vu_{\mathrm{int}} + \mathrm{dt}\,(\ve - \alpha \vu),
\qquad \vu = \vu_{\mathrm{int}} + k_p \ve, \label{eq:controller}\\
\vv_i &\leftarrow \vv_i + \tau ( \mW_i \vr_{i-1} + \vb_i + \mJ_i^\top \vu - \vv_i ),
\qquad \vr_i = \phi(\vv_i). \label{eq:dynamics}
\end{align}
We use $\mathrm{dt} = 0.02$, $\tau = \mathrm{dt}/0.2 = 0.1$, where $0.2$ is the ratio of the neuronal to the controller time constant, $k_p = 2.0$ and $\alpha = 1.7\times10^{-3}$, allow at most $500$ controller iterations, and stop when the change in the control signal falls below $10^{-4}$. The control strength is $\kappa = 0.3$ for modular arithmetic and $0.5$ for parity. The modular-arithmetic setting was selected on a grid over weight decay $\{10^{-3}, 3\times10^{-3}, 10^{-2}\}$ and $\kappa \in \{0.1, 0.2, 0.3\}$. At equilibrium each layer applies the local update
\begin{equation}
\Delta \mW_i \;\propto\; (\vr_i - \vr_i^{\mathrm{ff}})\, \vr_{i-1}^\top,
\label{eq:update}
\end{equation}
the difference between the controlled activity $\vr_i$ and the feedforward activity $\vr_i^{\mathrm{ff}}$, multiplied by the presynaptic activity. Its negative, $-2\,(\vr_i - \vr_i^{\mathrm{ff}})\,\vr_{i-1}^\top$ averaged over the batch, is written into the gradient slot, so that Adam's descent step moves the weights along $\Delta \mW_i$; it is clipped and passed to Adam like the other learning signals in this paper. Feedback through the transpose Jacobian follows the reference implementation of \citet{meulemans2021dfc}; the feedback pathway is therefore computed from the forward network. The error that enters the controller is a scaled squared-error gradient at the output, so DFC propagates a filtered loss gradient backward through the forward network, and its update differs from the loss gradient. DFC shows the same delayed generalization as the reference (Figure~\ref{fig:grok}). Evaluated every $10$ epochs, training accuracy reaches $1.0$ at epochs $90$ to $130$ for DFC and at epoch $100$ for the reference, and test accuracy first exceeds $0.95$ at epochs $650$ to $710$ for DFC and $1220$ to $1450$ for the reference; the grokking delays are $520$, $600$ and $610$ epochs for DFC and $1120$ to $1350$ for the reference over the three seeds, all above the $500$-epoch criterion of Section~\ref{sec:setup}.

\begin{figure}[!ht]
\begin{center}\includegraphics[width=0.5\linewidth]{figures/fig1_grok.pdf}\end{center}
\caption{Training (dashed) and test (solid) accuracy for the reference and DFC on modular addition ($p=31$), three seeds with shared initializations, evaluated every $10$ epochs; bands show one standard deviation across seeds, and the DFC runs end at epoch $1000$. Both fit the training set first and generalize hundreds of epochs later.}
\label{fig:grok}
\end{figure}

\paragraph{Measuring credit participation.}
\label{app:credit_measure}
Credit participation applies Equation~\ref{eq:pr} to the row norms of a neuron-level learning signal for the first hidden layer, one row per neuron, computed at every epoch of the last $800$ epochs and then averaged; it is the mean of instantaneous participation ratios. Which signal is logged depends on the condition (Table~\ref{tab:creditconv}). For the rules that hand a gradient-like signal to Adam it is that signal before clipping; for the adaptive optimizers, whose rescaling cannot be separated from their running state, it is the update applied to the first-layer weights. Circuit participation is measured from the outgoing weights of the trained network, on different parameters, so the agreement between credit and circuit participation is an empirical result and not an identity. Table~\ref{tab:baselines} also reports participation measured from Adam's effective step $m/\sqrt{v}$.

\begin{table}[!ht]
\caption{Signal on which credit participation is measured, per condition. All signals are for the first hidden layer, with one row per neuron.}
\label{tab:creditconv}
\begin{center}
\small
\begin{tabular}{@{}p{0.30\linewidth}p{0.66\linewidth}@{}}
\toprule
Condition & Logged signal \\
\midrule
Reference; LARS & raw loss gradient, before clipping and without the weight-decay term; for LARS this differs from the applied update only by the common layer factor and the weight-decay term \\
Credit equalization; group-size sweeps; imposed profiles & rescaled gradient of Equation~\ref{eq:knob} or~\ref{eq:grain}, before clipping and before Adam \\
DFC & local controller update of Equation~\ref{eq:update} for the first layer, before Adam \\
RMSProp, LAMB, Adafactor, Lion, Muon, Shampoo & update applied to the first-layer weights: gradient plus weight-decay term, transformed with the optimizer's running moments (Lion: sign of the momentum-interpolated gradient; LAMB: including the trust ratio) \\
\bottomrule
\end{tabular}
\end{center}
\end{table}
